\subsection{Coherent states and long-wavelength limit}
\label{sec:coherent-long-wavelength}

We next describe the common kinematical ingredients used to pass from the
spin chain to a smooth classical spin field.  We use standard $SU(2)$ spin
coherent states \cite{Perelomov1986}; their use in long-wavelength limits of
integrable spin chains is reviewed, for example, in
\cite{AvanDoikouSfetsos2010}.  The same construction is used in the direct
continuum analysis of the Class 5 and Class 6 chains
\cite{deLeeuwFontanellaNietoGarcia2026}.

Let $\sigma^1,\sigma^2,\sigma^3$ denote the Pauli matrices on the reference
local Hilbert space $\mathcal H_{\rm loc}\simeq\mathbb C^2$, and let
$|\!\uparrow\rangle$ and $|\!\downarrow\rangle$ be the eigenvectors of
$\sigma^3$ with eigenvalues $+1$ and $-1$, respectively.  The spin-$\tfrac12$
operators are normalized as
\begin{equation}
s^i:=\frac12\sigma^i,
\qquad i=1,2,3.
\label{eq:spin-half-normalization}
\end{equation}
Throughout, repeated Cartesian indices $i,j,k\in\{1,2,3\}$ are summed.
At the $n$th lattice site, a normalized spin coherent state is parametrized by
polar and azimuthal angles $\theta_n$ and $\phi_n$ as
\begin{equation}
|\Omega_n\rangle
=
\cos\frac{\theta_n}{2}\,|\!\uparrow\rangle
+e^{\ii\phi_n}\sin\frac{\theta_n}{2}\,|\!\downarrow\rangle .
\label{eq:spin-coherent-state}
\end{equation}
This phase convention fixes the associated unit vector.  We define the
\emph{classical spin vector at site $n$} by
\begin{equation}
\boldsymbol S_n
:=
\langle\Omega_n|\boldsymbol\sigma|\Omega_n\rangle
=
\begin{pmatrix}
\sin\theta_n\cos\phi_n\\
\sin\theta_n\sin\phi_n\\
\cos\theta_n
\end{pmatrix},
\qquad
\boldsymbol\sigma:=(\sigma^1,\sigma^2,\sigma^3)^{\mathsf T}.
\label{eq:site-classical-spin}
\end{equation}
It follows directly that
\begin{equation}
\boldsymbol S_n^{\mathsf T}\boldsymbol S_n=1,
\qquad
\langle\Omega_n|s^i|\Omega_n\rangle
=\frac12 S_n^i.
\label{eq:spin-coherent-expectation}
\end{equation}
Thus the normalization of the classical spin field is fixed independently of
the model-dependent normalization of the Lax matrix.

For a configuration of the full chain we use the normalized product coherent
state
\begin{equation}
|\Omega\rangle
:=
|\Omega_1\rangle\otimes\cdots\otimes|\Omega_N\rangle
\in\mathcal H^{(N)}.
\label{eq:product-coherent-state}
\end{equation}
The coherent-state labels are allowed to depend on time.  We denote by $t$ the
time coordinate of the continuum field theory.  Its relation to the quantum
lattice time $t_{\rm lat}$ introduced in \eqref{eq:Heisenberg-lattice-time}
is part of the continuum scaling.

Let $\epsilon>0$ denote the \emph{lattice spacing}.  We assign the spatial
coordinate
\begin{equation}
x_n:=(n-1)\epsilon,
\qquad n=1,\ldots,N,
\label{eq:lattice-coordinate}
\end{equation}
and take the long-wavelength limit
\begin{equation}
\epsilon\to0,
\qquad
N\to\infty,
\qquad
\ell:=N\epsilon\quad\text{fixed},
\label{eq:long-wavelength-limit}
\end{equation}
where $\ell$ is the physical circumference of the continuum spatial circle.
The long-wavelength assumption is that the coherent-state labels vary
smoothly from one lattice site to the next.  Equivalently, there is a smooth
unit spin field $\boldsymbol S(x,t)$ such that
\begin{equation}
\boldsymbol S_n(t)=\boldsymbol S(x_n,t),
\qquad
\boldsymbol S(x+\ell,t)=\boldsymbol S(x,t).
\label{eq:continuum-spin-field}
\end{equation}
For neighbouring sites this gives the Taylor expansion
\begin{equation}
\boldsymbol S_{n+1}
=
\boldsymbol S
+\epsilon\,\partial_x\boldsymbol S
+\frac{\epsilon^2}{2}\,\partial_x^2\boldsymbol S
+\mathcal O(\epsilon^3),
\label{eq:spin-long-wavelength-expansion}
\end{equation}
where the fields on the right-hand side are evaluated at $(x_n,t)$.
Differentiating the unit-spin constraint
$\boldsymbol S^{\mathsf T}\boldsymbol S=1$ gives the kinematical identities
\begin{equation}
\boldsymbol S^{\mathsf T}\partial_x\boldsymbol S=0,
\qquad
\boldsymbol S^{\mathsf T}\partial_x^2\boldsymbol S
=-(\partial_x\boldsymbol S)^{\mathsf T}(\partial_x\boldsymbol S).
\label{eq:spin-kinematical-identities}
\end{equation}
For any smooth scalar density $f(x)$, the corresponding lattice sum obeys the
usual Riemann-sum relation
\begin{equation}
\epsilon\sum_{n=1}^{N}f(x_n)
\longrightarrow
\int_0^{\ell}dx\,f(x).
\label{eq:riemann-sum-continuum}
\end{equation}
The powers of $\epsilon$ that accompany the Hamiltonian, deformation
parameters, and time coordinate are model dependent.

We now define the operation that maps local quantum operators to classical
quantities.  Let $X_a$ be an operator acting on the auxiliary space
$\mathcal V_a$ and on any finite set of physical lattice sites.  Its
\emph{coherent-state lower symbol}, denoted by the superscript $\downarrow$
\cite{BiskupChayesStarr2007}, is the auxiliary-space operator
\begin{equation}
X_a^{\downarrow}
:=
\bigl(\Id_{\mathcal V_a}\otimes\langle\Omega|\bigr)
X_a
\bigl(\Id_{\mathcal V_a}\otimes|\Omega\rangle\bigr).
\label{eq:lower-symbol-definition}
\end{equation}
Only the physical sites on which $X_a$ acts contribute to the matrix element;
the remaining factors in \eqref{eq:product-coherent-state} cancel by
normalization.  If $X$ has no auxiliary-space indices, $X^{\downarrow}$ is a
scalar function of the coherent-state labels.  In particular, the lower
symbol of the lattice Lax operator is
\begin{equation}
L_{a,n}^{\downarrow}(u)
=
\langle\Omega_n|L_{a,n}(u)|\Omega_n\rangle,
\label{eq:lattice-Lax-lower-symbol}
\end{equation}
which is a $2\times2$ matrix acting on the auxiliary space.

After evaluating these symbols on real spin vectors, we extend their
algebraic expressions to complex spin components satisfying
$\boldsymbol S^{\mathsf T}\boldsymbol S=1$.  Here $\mathsf T$ denotes the
transpose.  Scalar and vector products are bilinear on this complexified
spin space.

The continuum spatial Lax matrix is extracted after the model-dependent
normalization and spectral scaling have been fixed.  Let $\lambda$ denote the
\emph{continuum Lax spectral parameter}, to distinguish it from the quantum
spectral parameter $u$ defined in Section~\ref{sec:quantum-lattice-setup}.
For each model we choose a nonzero scalar normalization
$c_L(\epsilon,\lambda)$ and a relation $u=u(\epsilon,\lambda)$ such that the
normalized lattice operator
\begin{equation}
\widehat L_{a,n}(\epsilon,\lambda)
:=
c_L(\epsilon,\lambda)
L_{a,n}\bigl(u(\epsilon,\lambda)\bigr)
\label{eq:normalized-lattice-Lax}
\end{equation}
has a regular long-wavelength expansion.  When its lower symbol takes the
form
\begin{equation}
\widehat L_{a,n}^{\downarrow}(\epsilon,\lambda)
=
\Id_{\mathcal V_a}
+\epsilon\,U(x_n,t;\lambda)
+\mathcal O(\epsilon^2),
\label{eq:continuum-spatial-Lax-definition}
\end{equation}
we call the coefficient $U(x,t;\lambda)$ the \emph{continuum spatial Lax
matrix}.  The choices of $c_L$, of the scaling $u(\epsilon,\lambda)$, and of
any deformation parameters are part of the model-specific continuum limit;
they are specified separately for the two models.

The discrete zero-curvature equation \eqref{eq:discrete-zero-curvature}
also contains products of operators with overlapping physical support.
We next evaluate their lower symbols and determine their contribution to
the continuum zero-curvature equation.
